%=========================================================


\chapter*{Introducción}

El presente documento es el reporte técnico final del Trabajo Terminal con número de registro 2025-A138, titulado:

"Prototipo de sistema de identificación de estudiantes basado en reconocimiento facial dentro de la Escuela Superior de Cómputo para evitar la suplantación de identidad durante la aplicación de exámenes a título de suficiencia mediante el uso de credenciales y dispositivos móviles", elaborado por Alejandra De la Cruz De la Cruz, Luis Antonio Flores Esquivel, José Alfredo Jiménez Rodríguez y Daniel Martín Huertas Ramírez para obtener el grado de Ingeniero en Inteligencia Artificial.

Este trabajo fue dirigido por los profesores José Félix Serrano Talamantes y Ulises Vélez Saldaña, quienes fungieron como directores del proyecto.

El presente documento va dirigido a los sinodales, Andrés García Floriano,  Chadwick Carreto Arellano y Saúl de la O Torres, con la finalidad de exponer de manera detallada los objetivos, el desarrollo, los resultados y conclusiones del trabajo realizado, así como sustentar su validez técnica y académica para la evaluación correspondiente.

En este documento se describe el diseño, implementación,  pruebas y resultados de un sistema de identificación mediante reconocimiento facial, enfocado en asegurar la identidad de los estudiantes que presentan Exámenes a Titulo de Suficiencia en la Escuela Superior de Cómputo perteneciente al Instituto Politécnico Nacional.


% !TeX root = ../ejemplo2.tex

\section*{Organización del documento}

El presente documento se encuentra estructurado por los siguientes capítulos y anexos:

En el capítulo 1, se encuentran los antecedentes que permitieron identificar la problemática que rige este proyecto, lo que a su vez nos guío al análisis de esta misma para sustentar la propuesta de solución planteada a lo largo de este documento \\

En el capítulo 2, se presenta la definición del proyecto, donde se explica la justificación desde un contexto social, personal y académico del proyecto y además, se da una explicación breve y concisa sobre el alcance, los requerimientos, la arquitectura, las propiedades de software, el modelo de información y la metodología de desarrollo. \\

En el capítulo 3, se encuentra el estado del arte, en donde se revisaron y analizaron trabajos anteriores con similitudes al proyecto a realizar, esto con el fin de realizar una comparación de las principales características e identificar posibles áreas de mejora.\\

En el capítulo 4, se presenta el marco teórico donde se explican los términos, descripciones y demás elementos importantes que se relacionan en menor o mayor medida con este trabajo a fin de tener un panorama general sobre la terminología utilizada.\\

En el capítulo 5, se presenta los detalles del trabajo realizado, es decir se explica cómo y con que tecnologías y métodos se implementó cada una de las secciones del sistema.\\

En el capítulo 6, se exponen las pruebas realizadas en la aplicación móvil a los casos de uso más importantes, para buscar errores para su posterior corrección y/o determinar si la funcionalidad y/o proceso se está realizando de manera satisfactoria o cumple con lo esperado, además, se muestran las pruebas hechas al sistema de reconocimiento facial.\\

En el capítulo 7, se presentan los resultados obtenidos de las encuestas hechas a los alumnos, docentes y personal de seguridad de la escuela superior de cómputo (ESCOM), sobre el funcionamiento de la aplicación, las opiniones sobre esta y si les parece útil o si creen que podría ayudar reducir el problema de la suplantación de identidad durante los ETS. También, se presentan los resultados de las pruebas de la red neuronal y sus matrices de confusión.\\

Finalmente, en el capítulo 8, se presentan las direcciones futuras y consideraciones identificadas para este proyecto, una vez culminada la fase de Trabajo Terminal actual. Si bien se ha logrado la implementación de la aplicación móvil y el sistema de gestión asociado, incluyendo el reconocimiento facial en tiempo real, este capítulo delinea las potenciales expansiones y mejoras que podrían llevarse a cabo en una eventual continuación del proyecto, más allá de los objetivos académicos de este Trabajo Terminal.\\

En el anexo A, se sustenta la metodología de trabajo empleada, se presentan los resultados del estudio de factibilidad realizado y el análisis del sistema que compone la propuesta de solución a fin de delimitar el alcance de este y los elementos a desarrollar, asimismo, se presenta el análisis de riesgos correspondiente.\\  

En el Anexo B, se muestran las decisiones de diseño tomadas para este proyecto, esto se encuentra representado a través de diagramas para ofrecer una visualización y explicación del funcionamiento del sistema y las interacciones que existen entre los componentes que lo conforman. Estos incluyen el diagrama de estructura general del sistema, la estructura interna de los módulos, diagramas de clases, diagramas de componentes y diagramas de secuencia.\\



% Finalmente en el capítulo 6 se encuentra en trabajo futuro identificado para este trabajo terminal 1, destacando la implementación de la aplicación móvil, el sistema para la gestión de alumnos, personal y exámenes ETS así como la implementación del sistema de reconocimiento facial en tiempo real. Se espera que este trabajo se desarrolle y culmine durante la segunda etapa del trabajo terminal.

\chapter{Planteamiento del problema}

En este capítulo se describen los antecedentes y se delimita la problemática principal que guía el proyecto a realizar, así como los objetivos propuestos y la justificación de su desarrollo.

%---------------------------------------------------------
\input{introduccion/antecedentes}
%---------------------------------------------------------
\input{introduccion/planteamiento.tex}
%---------------------------------------------------------
\input{introduccion/propuesta.tex}

\chapter{Definición del proyecto}
Retomando lo dicho en la propuesta de solución, se llevó a cabo una aplicación móvil que sirve como herramienta para el personal de seguridad y el personal académico en temporadas de ETS para disminuir la suplantación de identidad dentro de la ESCOM. 

La aplicación móvil cuenta con varios apartados dependiendo del usuario que la esté usando, ya sea un alumno, un personal académico o un personal de seguridad, sin embargo, la parte fundamental de esta aplicación es la del docente y la del alumno en la que pueden realizar el reconocimiento facial. El alumno puede hacer el reconocimiento facial a su persona para rectificar que no haya ningún problema con su pase de asistencia al momento de llegar a realizar su ETS. El personal académico cuenta con la parte de reconocimiento facial como una forma de pasarle la asistencia al ETS al alumno. 

Es importante mencionar que el pase de asistencia no se realiza únicamente con el reconocimiento facial, se puede realizar mediante otras formas, como con un escanéo a la credencial del alumno o más fácil, si el personal académico ya conoce al alumno de antemano, lo puede especificar así al momento de registrar la asistencia al alumno. 

Por parte del personal de seguridad, puede registrar el ingreso a la unidad académica (ESCOM), por medio del escaneo de la credencial del alumno. 

La aplicación móvil cuenta con más secciones, tales como la visualización de los datos de los ETS, el calendario escolar del IPN, una sección de mensajes para que se los alumnos se puedan comunicar con el personal académico y para que el personal académico se pueda comunicar tanto con los alumnos como con otro personal académico. 

%---------------------------------------------------------
\input{introduccion/objetivos.tex}
%---------------------------------------------------------
\input{introduccion/justificación.tex}
%---------------------------------------------------------
\section{Alcance}
Este proyecto tiene como objetivo realizarse dentro de la Ciudad de México, específicamente dentro de la ESCOM del Instituto Politécnico Nacional. 

Específicamente, está pensado para que tanto el personal académico (específicamente para los docentes), el personal de seguridad y los alumnos, puedan usar la aplicación móvil planteada dentro de los periodos de ETS, tanto ETS ordinario como especial.

\section{Requerimientos}
El análisis de requerimientos arrojó un total de 26 requisitos de usuario (RU), distribuidos entre ambas plataformas: 17 para la aplicación móvil y 9 para el sistema web. Estos requerimientos fundamentales sirvieron como base para el diseño del sistema, generando 17 casos de uso para la plataforma móvil y 23 para la versión web.

Para consultar la documentación completa relacionada con este proceso, se recomienda revisar los siguientes apartados del Apéndice A:

\begin{itemize}
	\item \hyperref[cap:reqUsr]{\textbf{A.1 Modelo del Alcance}}: Contiene el detalle completo de los requerimientos de usuario identificados.
	\item \hyperref[cap:modDinamico]{\textbf{A.4 Modelo Dinámico}}: Incluye la especificación completa de casos de uso para ambas plataformas.
\end{itemize}


\section{Arquitectura}
El desarrollo de la arquitectura general de todo el proyecto presenta una arquitectura de servicios con una comunicación REST. Apesar de dividir ciertas tareas en diferentes servicios, se presenta una arquitectura monolítica de servicio, puesto que mayor parte de la lógica de negocio se encuentra dentro del servidor Java Spring Boot.

Por otra parte, el servidor que proporciona los servicios de reconocimiento facial se encuentra funcionando como API de forma aislada usando una arquitectura en capas compuesta por una capa de presentación que contiene los puntos de entrada de la API, una capa de aplicación o dominio que contiene las reglas de negocio de la aplicación y la funcionalidad del flujo de trabajo del sistema de reconocimiento facial y finalmente, una capa de persistencia que contiene la lógica para consultar y modificar las fuentes de datos necesarias.

Para una descripción más detallada sobre las estructuras internas de cada parte del sistema, así como la organización de cada sección, se puede consultar la anexo \ref{ApendiceB}


\section{Modelo de información}
Para realizar la parte del modelo de información se realizó un análisis que iba desde cubrir los datos de los alumnos, del personal académico y del personal de seguridad, así como las unidades académicas en las que estos trabajaban. 
Dentro de las unidades académicas también se vio que tuvieran programas académicos y dentro de estos, unidedades de aprendizaje.

Todo el análisis que se llevó a cabo se pensó para cubrir todo el proceso realizado desde que se inscribe un alumno a alguna unidad académica dentro del Instituto Politécnico Nacional, hasta que un alumno reprueba alguna unidad de aprendizaje. 

En el anexo \ref{sec:modeloDeDominioP} se detallan todas las entidades que se abstrajeron del análisis de todo este proceso, a su vez, se detalla cada una de las variables que estas entidades contienen.

\section{Metodología de desarollo}
Se adoptó el modelo en espiral como metodología principal, dado su enfoque iterativo y adaptable que permite incorporar mejoras progresivas, gestionar riesgos y ajustar los requerimientos durante cada fase. Este modelo combina la estructuración del enfoque en cascada con la flexibilidad para evolucionar conforme a las necesidades del proyecto.

El proceso se organizó en seis iteraciones clave, desde la investigación inicial hasta la implementación de prototipos funcionales. En cada ciclo se definieron objetivos específicos, se evaluaron riesgos y se generaron productos intermedios para garantizar el avance controlado del sistema. Para una descripción detallada de la metodología aplicada, consúltese el \hyperref[cap:Meto]{\textbf{Anexo A.3 Metodología}}.
%---------------------------------------------------------